\section{Conclusion}\label{sec:conclusion}

We have shown that the implication structure of $4{,}694$ equational laws on
magmas, despite arising from a computationally intensive formal verification
effort, admits a surprisingly compact description.  A five-class structural
form classification resolves nearly 40\% of all implications with zero error,
and a small set of additional rules based on equation signatures and variable
counts brings the total accuracy to 88.9\% on random samples.

The key insight is hierarchical: absorbing laws universally imply all other
laws (Theorem~\ref{thm:absorbing}), general-form laws never imply standard-form
laws (Theorem~\ref{thm:general-nonimplication}), signature direction constrains
the flow of implications (Theorem~\ref{thm:sig-direction}), and variable count
difference provides a calibrated probability for the remaining cases
(Proposition~\ref{prop:nvar-monotone}).  Each layer addresses a distinct subset
of the implication space with decreasing certainty but increasing coverage.

Our work contributes to the broader question of whether mathematical knowledge
can be distilled into compact, interpretable representations.  The equational
theories domain provides a clean test case: the ground truth is complete, the
objects are combinatorially structured, and the patterns are empirically
verifiable.  The success of simple structural rules suggests that similar
distillation may be possible in other areas of algebra where implication
lattices exhibit regularities.

Future directions include extending the analysis to higher-order laws,
incorporating Stone pairing statistics for improved prediction in the ``hard
middle'' cases, and developing theoretical bounds on the achievable accuracy
as a function of rule complexity.
